\section{Criteria for infinite finitistic dimension}
\label{sec:criteria}

Both constructions use a module with prescribed homological behaviour. The
following dimension-shifting argument produces modules of every finite
projective dimension.

\begin{proposition}
  \label{prop:coresolution}
  Let $\A$ be an abelian category with enough projective objects, and suppose
  given short exact sequences
  \[
    0\to c_0\to p_0\to c_1\to0\qquad\text{and}\qquad
    0\to c_n\to p_n\to c_{n+1}\to0\quad(n\geq1)
  \]
  in which every $p_n$ is projective, $c_0$ is projective and $c_1$ is not.
  Then ${\operatorname{pd}c_n=n}$ for every ${n\geq1}$.
\end{proposition}

\begin{proof}
  A short exact sequence ${0\to c_n\to p_n\to c_{n+1}\to0}$ with $p_n$
  projective gives ${\operatorname{pd}c_{n+1}\leq\operatorname{pd}c_n+1}$, so
  ${\operatorname{pd}c_n\leq n}$ by induction. Since $c_1$ is not projective,
  ${\operatorname{pd}c_1=1}$. Let ${n\geq1}$ and suppose that
  ${\operatorname{pd}c_n=n}$. For every object $V$ and every ${j\geq1}$, the long
  exact sequence of ${\operatorname{Ext}(-,V)}$ and the vanishing of
  $\operatorname{Ext}^{\geq1}(p_n,V)$ give an isomorphism
  ${\operatorname{Ext}^{j+1}(c_{n+1},V)\cong\operatorname{Ext}^j(c_n,V)}$.
  Choosing $V$ with ${\operatorname{Ext}^n(c_n,V)\neq0}$ and ${j=n}$ yields
  ${\operatorname{pd}c_{n+1}=n+1}$.
\end{proof}

\begin{remark}
  \Cref{prop:coresolution} is formally verified in Lean, against Mathlib, in an
  arbitrary abelian category; see \Cref{app:verification}.
\end{remark}

\subsection{The strong Nakayama conjecture}
\label{subsec:strong-nakayama}

The following argument shows that finiteness of the finitistic dimension of
$A^{\op}$ implies the strong Nakayama conjecture for $A$~\cite[Section~3.2,
Proposition~5]{CB19}; it also gives the exact projective dimensions of the
modules constructed.

\begin{theorem}
  \label{thm:strong-nakayama}
  Let $A$ be a finite-dimensional algebra and ${E\neq0}$ a right $A$-module
  such that
  \[
    \operatorname{Ext}^i_{A^{\op}}(E,A)=0\qquad\text{for all }i\geq0.
  \]
  Let ${\cdots\to P_1\to P_0\to E\to0}$ be a minimal projective resolution, and
  put ${C_n}={\operatorname{coker}(P_{n-1}^*}\to{ P_n^*)}$ for ${n\geq1}$. Then
  ${C_n\cong\operatorname{Tr}\Omega^{n-1}E}$ and ${\operatorname{pd}_AC_n=n}$ for
  every ${n\geq1}$. In particular, ${\operatorname{findim}A=\infty}$.
\end{theorem}

\begin{proof}
  The complex ${0\to P_0^*\to P_1^*\to P_2^*\to\cdots}$ of projective left
  $A$-modules has cohomology ${\operatorname{Ext}^*_{A^{\op}}(E,A)=0}$, so it is
  exact, and it splits into short exact sequences
  ${0\to C_0\to P_1^*\to C_1\to0}$ and ${0\to C_n\to P_{n+1}^*\to C_{n+1}\to0}$,
  where ${C_0=P_0^*}$. If $C_1$ were projective, the map ${P_0^*\to P_1^*}$ would be
  a split monomorphism; since finitely generated projective modules are
  reflexive, its dual ${P_1\to P_0}$ would then be a split epimorphism, and
  ${E=\operatorname{coker}(P_1\to P_0)}$ would vanish. Hence,
  \Cref{prop:coresolution} applies. Finally,
  ${P_n\to P_{n-1}\to\Omega^{n-1}E\to0}$ is a minimal presentation, so that
  \[
    \operatorname{Tr}\Omega^{n-1}E=\operatorname{coker}(P_{n-1}^*\to P_n^*)
    =C_n.\qedhere
  \]
\end{proof}

Since ${\operatorname{Ext}^i_{A^{\op}}(E,A)\cong D\operatorname{Tor}^A_i(E,DA)}$,
the hypothesis says that the derived Nakayama functor ${-\Lotimes[A]DA}$
annihilates $E$. It forces ${\operatorname{pd}E=\infty}$: if
${\operatorname{pd}E=n<\infty}$, then ${\operatorname{Ext}^n(E,A)\neq0}$.

\begin{remark}
  \Cref{thm:strong-nakayama} is formally verified in Lean over an arbitrary
  ring, including the equivalence between the vanishing of the groups
  $\operatorname{Ext}^i_{A^{\op}}(E,A)$ and the exactness of the dual complex;
  the identification of $C_n$ with a transpose is proved for the presentation
  given by the resolution, without its minimality. See
  \Cref{app:verification}.
\end{remark}

The next proposition reformulates the hypothesis of \Cref{thm:strong-nakayama}
as a property of a module of projective dimension one, such as $C_1$.

\begin{proposition}
  \label{prop:infinitely-torsionfree}
  Let $C$ be an $A$-module without nonzero projective summands and with
  ${\operatorname{pd}C=1}$. The following are equivalent:
  \begin{enumerate}
    \item the right module ${E=\operatorname{Tr}C}$ satisfies
      ${\operatorname{Ext}^i_{A^{\op}}(E,A)=0}$ for all ${i\geq0}$;
    \item $C$ is not projective and
      ${\operatorname{Ext}^i_{A^{\op}}(\operatorname{Tr}C,A)=0}$ for all ${i\geq1}$;
    \item $C$ is not projective and admits an exact sequence
      ${0\to C\to Q_2\to Q_3\to\cdots}$ with each $Q_m$ projective, in which every
      map from $C$ or from a cokernel to the next $Q_m$ is a left
      $\operatorname{add}A$-approximation.
  \end{enumerate}
  Moreover, ${E\cong\operatorname{Ext}^1_A(C,A)}$.
\end{proposition}

\begin{proof}
  Choose a minimal projective resolution ${0\to Q_0\to Q_1\to C\to0}$. Then
  ${E=\operatorname{coker}(Q_1^*\to Q_0^*)\cong\operatorname{Ext}^1_A(C,A)}$ and
  ${E^*=\ker(Q_0\to Q_1)=0}$; moreover ${E\neq0}$, since otherwise ${Q_1^*\to Q_0^*}$
  would split and $C$ would be projective. This proves that (1) and (2) are
  equivalent. Assume (1). Dualising a minimal projective resolution of $E$ gives
  an exact complex whose first cokernel is ${\operatorname{Tr}E=C}$, since
  ${\operatorname{Tr}\operatorname{Tr}C\cong C}$ for a module $C$ without nonzero
  projective summands, and whose tail
  is a projective coresolution of $C$; the dual of each of its short exact
  sequences is a segment of the resolution of $E$, hence exact, which is the
  approximation property in (3). Conversely, dualising the sequence in (3) gives
  an exact sequence ${\cdots\to Q_3^*\to Q_2^*\to C^*\to0}$; splicing it with
  ${0\to C^*\to Q_1^*\to Q_0^*\to E\to0}$ gives a projective resolution of $E$
  whose dual ${0\to Q_0\to Q_1\to Q_2\to\cdots}$ is exact, and (1) follows.
\end{proof}

\begin{corollary}
  \label{coro:infinitely-many-torsionless}
  If a finite-dimensional algebra $A$ admits a module $E$ as in
  \Cref{thm:strong-nakayama}, then $A$ has infinitely many isomorphism classes
  of indecomposable torsionless modules of finite projective dimension.
\end{corollary}

\begin{proof}
  The modules $C_n$ of \Cref{thm:strong-nakayama} embed into the projective
  modules $P_{n+1}^*$, by the short exact sequences in its proof, so their
  indecomposable summands are torsionless and of finite projective
  dimension. If there were only finitely many of them, their projective
  dimensions would bound ${\operatorname{pd}C_n=n}$ for all $n$.
\end{proof}

\subsection{Counterexamples to the Auslander--Reiten conjecture}
\label{subsec:auslander-reiten}

Let $\Lambda$ be an Artin algebra. The \emph{Auslander--Reiten conjecture}
asserts that a $\Lambda$-module $M$ with
${\operatorname{Ext}^i_\Lambda(M,M\oplus\Lambda)=0}$ for all ${i\geq1}$ is
projective. Auslander and Reiten showed that its failure
produces a failure of the generalised Nakayama conjecture for an endomorphism
algebra~\cite[Theorem~1.1(b)]{AR75a}. Combined with \Cref{thm:strong-nakayama},
this gives a counterexample to the finitistic dimension conjecture with one
more simple module.

\begin{theorem}
  \label{thm:ar-to-findim}
  Let $\Lambda$ be a finite-dimensional algebra and $M$ a nonprojective
  $\Lambda$-module with ${\operatorname{Ext}^i_\Lambda(M,M\oplus\Lambda)}={0}$ for
  all ${i\geq1}$. Put ${G=\Lambda\oplus M}$ and
  ${\Gamma}={\operatorname{End}_\Lambda(G)^{\op}}$,
  and let ${\pi\colon P\to M}$ be a projective cover. Then the $\Gamma$-module
  \[
    S\coloneqq\operatorname{coker}\bigl(\Hom[\Lambda]{G}{\pi}\bigr)
  \]
  is nonzero and ${\operatorname{Ext}^i_\Gamma(S,\Gamma)=0}$ for all ${i\geq0}$.
  Consequently, the algebra $\operatorname{End}_\Lambda(G)$ has infinite
  little finitistic dimension.
\end{theorem}

\begin{proof}
  Write ${\mathbb{F}=\Hom[\Lambda]{G}{-}}$, which restricts to an equivalence
  ${\operatorname{add}G\simeq\operatorname{proj}\Gamma}$ with
  ${\mathbb{F}G=\Gamma}$.
  Put ${K_0=\ker\pi}$ and choose right $\operatorname{add}G$-approximations
  ${g_j\colon G_j\to K_j}$ with ${K_{j+1}=\ker g_j}$; they are surjective since
  ${\Lambda\in\operatorname{add}G}$. The functor $\mathbb{F}$ is left exact and
  the
  maps $\mathbb{F}g_j$ are surjective, so
  \[
    \cdots\to\mathbb{F}G_1\to\mathbb{F}G_0\to\mathbb{F}P\to\mathbb{F}M\to S\to0
  \]
  is a projective resolution of $S$. If $\mathbb{F}\pi$ were surjective, the
  identity of the summand $M$ of $G$ would lift through $\pi$, and $M$ would be
  projective; hence ${S\neq0}$. For ${Y\in\operatorname{add}G}$ there is a natural
  isomorphism ${\Hom[\Gamma]{\mathbb{F}Y}{\Gamma}\cong\Hom[\Lambda]{Y}{G}}$, so
  that
  $\operatorname{Ext}^*_\Gamma(S,\Gamma)$ is the cohomology of
  \[
    0\to\Hom[\Lambda]{M}{G}\to\Hom[\Lambda]{P}{G}\to\Hom[\Lambda]{G_0}{G}\to
    \Hom[\Lambda]{G_1}{G}\to\cdots,
  \]
  with $\Hom[\Lambda]{M}{G}$ in degree zero. Since
  ${\operatorname{Ext}^{\geq1}_\Lambda(G,G)}={0}$, the resolution
  ${\cdots}\to{ G_0}\to{ P}\to{ M}$ consists of $\Hom[\Lambda]{-}{G}$-acyclic modules
  and
  computes $\operatorname{Ext}^*_\Lambda(M,G)$; as
  ${\operatorname{Ext}^{\geq1}_\Lambda(M,G)}={0}$ and the first map is injective,
  the complex is exact. The last assertion is \Cref{thm:strong-nakayama} applied
  to ${\Gamma^{\op}}={\operatorname{End}_\Lambda(G)}$.
\end{proof}

Replacing $M$ by an indecomposable nonprojective summand preserves the
hypotheses of \Cref{thm:ar-to-findim}, since $\operatorname{Ext}$ is additive.
If $\Lambda$ is basic, the algebra $\operatorname{End}_\Lambda(G)$ then has
exactly one simple module more than $\Lambda$.

Conversely, simple modules witnessing failure of the strong Nakayama conjecture
always arise in this
way, on a corner of the algebra. In the next proposition, $Af$, $M$ and their
endomorphisms are taken as \emph{right} $B$-modules.

\begin{proposition}
  \label{prop:simple-witness}
  Let $S$ be a simple right $A$-module with
  ${\Hom[A^{\op}]{S}{A}}={0}={\operatorname{Ext}^1_{A^{\op}}(S,A)}$. Let $e$ be the
  sum
  of the primitive idempotents of a complete decomposition of $1$ corresponding
  to $S$, ${f=1-e}$, ${B=fAf}$ and ${M=eAf}$. Then left multiplication induces an
  algebra isomorphism
  \[
    A\stackrel{\sim}{\longrightarrow}\operatorname{End}_{B^{\op}}(Af)
    =\operatorname{End}_{B^{\op}}(B\oplus M),
  \]
  $B$ has exactly one isomorphism class of simple modules fewer than $A$, and
  $M$ is not a projective right $B$-module. If moreover
  ${\operatorname{Ext}^i_{A^{\op}}(S,A)=0}$ for all ${i\geq0}$, then
  ${\operatorname{Ext}^i_{B^{\op}}(M,B\oplus M)=0}$ for all ${i\geq1}$.
\end{proposition}

\begin{proof}
  For a right $A$-module $V$, let ${\eta_V\colon V\to\Hom[B^{\op}]{Af}{Vf}}$ send
  $v$ to ${x\mapsto vx}$. Multiplication by $f$ identifies $\eta_Vf$ with the
  identity of $Vf$, so the kernel and cokernel of $\eta_V$ are modules over
  $A/AfA$, whose composition factors are copies of $S$. If
  ${\Hom[A^{\op}]{S}{V}=0}$, then $\eta_V$ is injective. If moreover
  ${\operatorname{Ext}^1_{A^{\op}}(S,V)=0}$, then
  ${\operatorname{Ext}^1_{A^{\op}}(T,V)=0}$ for every $A/AfA$-module $T$ by
  induction on the length, so $\eta_V$ splits; its cokernel is a summand of
  $\Hom[B^{\op}]{Af}{Vf}$ annihilated by $f$, hence zero, since
  ${\Hom[A^{\op}]{T}{\Hom[B^{\op}]{Af}{Vf}}\cong\Hom[B^{\op}]{Tf}{Vf}=0}$. For
  ${V=A}$ this gives the isomorphism. Passing from $A$ to $B$ removes exactly the
  isomorphism class of $S$. If $M$ were projective, ${B\oplus M}$ would be a
  progenerator of $B$ and $A$ would be Morita equivalent to $B$, which has fewer
  simple modules.

  Assume now the vanishing in all degrees. A minimal injective resolution
  ${A\to I^\bullet}$ of the right regular module has no summand with socle $S$,
  since the number of such summands of $I^j$ is the dimension of
  ${\operatorname{Ext}^j_{A^{\op}}(S,A)}$ over ${\operatorname{End}_{A^{\op}}(S)}$,
  so each $I^j$ is a sum of modules $D(Ae_t)$ with ${e_t\leq f}$. Then
  $I^jf$ is a direct sum of modules $D(Be_t)$, hence an injective right
  $B$-module, and, by the first part of the proof applied to ${V=I^j}$, the
  complex
  ${\Hom[B^{\op}]{Af}{I^\bullet f}\cong\Hom[A^{\op}]{A}{I^\bullet}}$ computes
  $\operatorname{Ext}^*_{B^{\op}}(Af,Af)$, which therefore vanishes in positive
  degrees.
\end{proof}

\begin{example}
  The two vanishing conditions in the first part of
  \Cref{prop:simple-witness} do not suffice for the last assertion. Let $S$ be
  the right simple module at the vertex $3$ of the algebra
  ${A}={\kk(1\xrightarrow{a}2\xrightarrow{b}3)/(ba)}$. Its minimal projective resolution is
  ${0\to e_1A\to e_2A\to e_3A\to S\to0}$, and the dual complex shows that
  ${\Hom[A^{\op}]{S}{A}}={0}={\operatorname{Ext}^1_{A^{\op}}(S,A)}$, while
  ${\operatorname{Ext}^2_{A^{\op}}(S,A)}\cong{ Ae_1/\kk a}\neq{0}$. The
  conclusion fails as well: here ${B=\kk(1\xrightarrow{a}2)}$ and $M$ is the simple
  right $B$-module at the vertex $2$, with projective resolution
  ${0\to e_1B\to e_2B\to M\to0}$, so that
  ${\operatorname{Ext}^1_{B^{\op}}(M,B)\cong Be_1/\kk a\neq0}$.
\end{example}

\subsection{Triangular matrix algebras and bounds on projective dimension}
\label{subsec:gluing}

Triangular matrix algebras cannot produce a counterexample to the strong
Nakayama conjecture when both
diagonal algebras have finite little finitistic dimension.

\begin{proposition}
  \label{prop:triangular}
  Let ${A=\begin{psmallmatrix}B&0\\M&C\end{psmallmatrix}}$ for a $C$-$B$-bimodule
  $M$, so that right $A$-modules are triples $(X,Y,\varphi)$ with $X$ a right
  $B$-module, $Y$ a right $C$-module and ${\varphi\colon Y\otimes_CM\to X}$. Let
  ${E=(X,Y,\varphi)\neq0}$ satisfy ${\RHom[A^{\op}]{E}{A}=0}$, and put
  ${Z=\operatorname{cone}(Y\Lotimes[C]M\to X)}$ in
  ${\DerCat[-]{\operatorname{mod}B^{\op}}}$.
  Then ${\RHom[B^{\op}]{Z}{B}=0}$. If ${Z\neq0}$, the little finitistic dimension of
  $B$ is infinite; if ${Z=0}$, then ${Y\neq0}$, ${\RHom[C^{\op}]{Y}{C}=0}$, and the
  little finitistic dimension of $C$ is infinite.
\end{proposition}

\begin{proof}
  Let ${e=\operatorname{diag}(1,0)}$ and ${f=\operatorname{diag}(0,1)}$, let $i_*$
  inflate right $B$-modules and put ${j_!Y=Y\Lotimes[C]fA}$. The derived counit
  gives
  a triangle ${j_!Y\to E\to i_*Z\to(j_!Y)[1]}$. Since ${eA=i_*B}$, adjunction gives
  ${\RHom[B^{\op}]{Z}{B}\cong\RHom[A^{\op}]{E}{eA}=0}$. If ${Z\neq0}$, represent it
  by a nonzero minimal bounded-above complex ${P^\bullet}$ of finitely generated
  projective right $B$-modules, with largest nonzero degree $m$. Its dual
  ${0\to(P^m)^*\to(P^{m-1})^*\to\cdots}$ is exact, and its first cokernel is not
  projective, since otherwise the first map would split, against the
  minimality of ${P^\bullet}$; \Cref{prop:coresolution} applied to the successive
  cokernels gives unbounded finite projective dimensions. If ${Z=0}$, then
  ${E\simeq j_!Y}$, so ${Y\neq0}$, and
  ${\RHom[C^{\op}]{Y}{C}\cong\RHom[A^{\op}]{j_!Y}{fA}=0}$ since ${(fA)f=C}$; now
  \Cref{thm:strong-nakayama} applies to $C$.
\end{proof}

Finite projective dimensions are bounded on modules of bounded
dimension. The argument is due to Schofield, as explained by Happel in his
survey~\cite[Section~2.3]{Hap90}; the openness also follows from the
semicontinuity of $\dim\operatorname{Ext}$ on module
varieties~\cite[Corollary~2.6]{GLS24}. Both sources work over an algebraically
closed field; the argument below works over any field.

\begin{proposition}
  \label{prop:bounded-dimension}
  Let $A$ be a finite-dimensional algebra and ${d\geq0}$ an integer. The
  subset of the variety $\operatorname{Rep}_d(A)$ of $d$-dimensional modules of
  projective dimension at most $n$ is open, and projective dimension is bounded
  on the modules of finite projective dimension in $\operatorname{Rep}_d(A)$.
  In particular, modules of unbounded finite projective dimension have unbounded
  dimension.
\end{proposition}

\begin{proof}
  Let $J$ be the radical of $A$. A module $M$ satisfies
  ${\operatorname{pd}M\leq n}$ if and only if
  ${\operatorname{Tor}^A_{n+1}(A/J,M)=0}$. Choose a resolution ${Q_\bullet}$ of the
  right module $A/J$ by finitely generated free modules. Over
  $\operatorname{Rep}_d(A)$, the complex
  \[
    Q_{n+2}\otimes_AM\longrightarrow Q_{n+1}\otimes_AM\longrightarrow
    Q_n\otimes_AM
  \]
  has terms of fixed dimension and differentials given by matrices whose entries
  are polynomial in the point $M$; the vanishing of its middle homology is a
  lower bound on the sum of the ranks of the two differentials, which is an open
  condition. The open subsets ${\set{M}[\operatorname{pd}M\leq n]}$ form an
  ascending chain in a noetherian topological space and are therefore
  stationary.
\end{proof}

In \Cref{thm:strong-nakayama}, the dimension of $\operatorname{Tr}N$ is at
most $(\dim A)^2\dim N$. Hence, by \Cref{prop:bounded-dimension}, the syzygies
$\Omega^nE$ of a module $E$ as in \Cref{thm:strong-nakayama} have unbounded
dimension:
families of modules of constant dimension whose syzygies move along a
parameter, such as those of Schulz~\cite[Example~7]{Schu95}, cannot themselves
satisfy the hypothesis of that theorem.
